% ID: SEL-D0015-BC-E06
% ESTADO: cotejada. Segunda pasada visual completa y borrador compilado revisados.
% Fuente: Selectividad/matematicas_aplicadas/matematicas_aplicadas_P_C.pdf
% SHA-256: 463bf63d4697a18a0b9efc9287e7b0a593c803bab98a3cfab0e260067a0dd02c
% Página: 2; bloque C; ejercicio original 6.

Las bicicletas de alquiler de una ciudad se clasifican por su calidad:
buena, media y mala. El 30\% de dichas bicicletas son gestionadas por una
empresa $E_1$ y el resto por una empresa $E_2$. De las bicicletas de la empresa
$E_1$, el 80\% son de buena calidad, el 5\% de calidad media y el resto de mala
calidad. De las bicicletas de la empresa $E_2$ se sabe que el 60\% son de buena
calidad, pero se desconocen los porcentajes de bicicletas de calidad media
y calidad mala. Se elige al azar una bicicleta de alquiler de esa ciudad.
\begin{enumerate}[label=\alph*)]
  \item \textbf{(1 punto)} Calcule la probabilidad de que sea de buena calidad.
  \item \textbf{(0.75 puntos)} Calcule la probabilidad de que sea de la empresa
  $E_1$ y de mala calidad.
  \item \textbf{(0.75 puntos)} Si se sabe que el porcentaje de bicicletas de
  alquiler de calidad media en toda la ciudad es del 19\%,
  \textquestiondown{}cuál es la probabilidad de que sea de calidad media,
  sabiendo que la bicicleta elegida es de la empresa $E_2$?
\end{enumerate}
